\documentclass[10pt]{article}
\usepackage{graphicx}
\usepackage{helvet}
\renewcommand{\familydefault}{\sfdefault}
\usepackage{hyperref}
\usepackage{enumitem}
\usepackage[utf8]{inputenc} 
\usepackage[T1]{fontenc}
\usepackage[english]{babel}
\usepackage[left=1cm,top=1cm,right=1cm,bottom=1cm]{geometry}
\usepackage{xcolor}

\definecolor{mylinkcolor}{RGB}{0,0,150}
\hypersetup{
	colorlinks=true,
	linkcolor=mylinkcolor,
	filecolor=mylinkcolor,      
	urlcolor=mylinkcolor,
	citecolor=mylinkcolor,
}

\setlength{\parindent}{0pt}
\setlength{\parskip}{2pt} 

\begin{document}
	\pagestyle{empty}
	\begin{center}
		\textbf{\Large Yiquan Wang}\\
		\vspace{3pt}
		+86-19537838515 \textbullet \ \href{mailto:ethan@stu.xju.edu.cn}{ethan@stu.xju.edu.cn} \textbullet \ \href{https://wyqmath.cn/}{Personal Homepage} \textbullet \ \href{https://github.com/wyqmath}{GitHub} \textbullet \ \href{https://scholar.google.com/citations?user=ULP3e1cAAAAJ}{Google Scholar} \textbullet \ \href{https://openreview.net/profile?id=~Yiquan_Wang3}{OpenReview} \textbullet \ \href{https://orcid.org/0000-0003-1417-5752}{ORCID} \\
		\vspace{2pt}
		\textit{Memberships:} Biophysical Society of China \textbullet \ Chinese Chemical Society \textbullet \ IEEE Biometrics Council \\
		\vspace{-2pt}
		\hrulefill
	\end{center}
	
	\vspace{-2pt}
	\begin{center}
		\textbf{\large Education}
	\end{center}
	\vspace{-4pt}
	\textbf{Xinjiang University (National Base for Mathematical Research and Teaching Talents)} \hfill Urumqi, Xinjiang \\
	\textit{Bachelor of Science in Mathematics and Applied Mathematics} \hfill 2023.09 – 2027.06 
	
	\vspace{2pt}
	\textbf{Tsinghua University \& X-Institute} \hfill Shenzhen, Guangdong \\
	\textit{Tsien Excellence in Engineering Program (Joint Program, X Scholar)} \hfill 2024.06 – 2027.06
	
	\vspace{2pt}
	\begin{center}
		\textbf{\large Research Experience}
	\end{center}
	\vspace{-4pt}
	\textbf{Beijing Frontier Research Center for Biological Structure, Tsinghua University} \hfill Beijing, China \\
	\textit{Research Intern | Supervised by \href{https://scholar.google.com/citations?user=gnJCO14AAAAJ}{Prof. Yafei Yuan} \& \href{https://wanggroup.ai/}{Prof. Tong Wang}} \hfill 2025.07 – Present
	\begin{itemize}[noitemsep, topsep=2pt, partopsep=0pt, parsep=0pt, leftmargin=*]
		\item Engaged in structure-based protein and polypeptide design research leveraging generative Artificial Intelligence.
		\item Investigated protein conformational changes and functional mechanisms through AI-driven molecular dynamics (MD) simulations.
		\item Independently developed high-throughput data infrastructures for the structural biology community: the \href{https://www.fbs.frcbs.tsinghua.edu.cn/competition/2025Peptide/}{Peptide Design Competition Website} and the \href{https://www.frcbs.tsinghua.edu.cn/cpdb/}{Polypeptide Structure Database}.
	\end{itemize}
	
	\vspace{2pt}
	\textbf{Shenzhen Bay Laboratory} \hfill Shenzhen, Guangdong \\
	\textit{Visiting Student | Institute of Neurological and Psychiatric Disorders (\href{https://www.szbl.ac.cn/research/group/Wen-Yuan.htm}{Wen Yuan's Group})} \hfill 2025.07 – 2025.09
	\begin{itemize}[noitemsep, topsep=2pt, partopsep=0pt, parsep=0pt, leftmargin=*]
		\item Engaged in interdisciplinary research across computational biology, neuroscience, and biochemistry, integrating high-throughput single-cell multi-omics and CRISPR screening to investigate the pathogenic mechanisms and therapeutic targets of neurological disorders (e.g., Alzheimer's and autism).
	\end{itemize}
	
	\vspace{2pt}
	\textbf{AI-Driven Protein Function Prediction \& Generative Design} (Tsinghua University ESRT) \hfill 2024.08 – 2025.10 \\
	\textit{Project Lead | Supervised by Prof. Kai Wei}
	\begin{itemize}[noitemsep, topsep=2pt, partopsep=0pt, parsep=0pt, leftmargin=*]
		\item Proposed protein sonification (1D to 2D spectrograms); built a multimodal fusion model achieving 90.44\% accuracy on CARE enzyme benchmark, matching standard transformers (ESM-2, ProtBERT) in data efficiency.
		\item Guided a Diffusion Model using the novel Two-Dimensional encoding to generate Green Fluorescent Protein variants, validated via biophysics.
	\end{itemize}
	
	\vspace{2pt}
	\begin{center}
		\textbf{\large Selected Publications}
	\end{center}
	\vspace{-4pt}
	\textbf{AI-Driven Protein Modeling, Design \& Structural Bioinformatics}
	\begin{enumerate}[noitemsep, topsep=2pt, partopsep=0pt, parsep=0pt, leftmargin=*, label={[\arabic*]}]
		\item \textbf{Wang, Y.}, Cai, M., Dong, Y., Ma, Y., \& Wei, K. (2025). \href{https://pubs.acs.org/doi/abs/10.1021/acs.jcim.5c01768}{From Signal to Symphony: Exploring 2D Sequence Representations for Protein Function Prediction}. \textit{Journal of Chemical Information and Modeling}, 65(23), 12723-12736. (JCR Q1, CAS Q2 \textbf{Top}).
		\item \textbf{Wang, Y.}, Cai, M., Ma, Y., \& Wei, K. (2026). \href{https://openreview.net/forum?id=rTO3JOgC3a}{Sequence-context-aware decoding enables robust reconstruction of protein dynamics from crystallographic B-factors}. Accepted to \textit{Learning Meaningful Representations of Life (LMRL) Workshop at ICLR 2026}, under review at \textit{Science China Life Sciences}.
		\item Wang, X.$^\dagger$, \textbf{Wang, Y.}$^\dagger$, Huang, T.-Y., Dong, Y., Deng, J., Xu, L., Li, X., \& He, R. (2025). \href{https://ieeexplore.ieee.org/abstract/document/11356718/}{Octopus Inspired Optimization (OIO): A Hierarchical Framework for Navigating Protein Fitness Landscapes}. \textit{2025 IEEE International Conference on Bioinformatics and Biomedicine (BIBM)}, 6897-6904. (CCF-B, Co-first author).
		\item \textbf{Wang, Y.}, Ma, Y., Chang, Y., Yan, J., Zhang, J., Cai, M., \& Wei, K. (2025). \href{https://www.mdpi.com/2079-7737/14/12/1665}{Diffusion Models at the Drug Discovery Frontier: A Review on Generating Small Molecules Versus Therapeutic Peptides}. \textit{Biology}, 14(12), 1665. (JCR Q1, CAS Q2).
	\end{enumerate}
	
	\vspace{2pt}
	\textbf{First-Principles of Protein Geometry via Mathematical Physics}
	\begin{enumerate}[noitemsep, topsep=2pt, partopsep=0pt, parsep=0pt, leftmargin=*, label={[\arabic*]}, resume]
		\item \textbf{Wang, Y.} (2026). \href{https://arxiv.org/abs/2602.21787}{Spectral entropy of the discrete Hasimoto effective potential exposes sub-residue geometric transitions in protein secondary structure}. \textit{arXiv preprint arXiv:2602.21787}.
		\item \textbf{Wang, Y.} (2026). \href{https://arxiv.org/abs/2602.16255}{Piecewise integrability of the discrete Hasimoto map for analytic prediction and design of helical peptides}. \textit{arXiv preprint arXiv:2602.16255}.
		\item \textbf{Wang, Y.} (2026). \href{https://arxiv.org/abs/2602.13160}{Structural barriers of the discrete Hasimoto map applied to protein backbone geometry}. \textit{arXiv preprint arXiv:2602.13160}.
	\end{enumerate}
	
	\vspace{2pt}
	\begin{center}
		\textbf{\large Selected Awards \& Academic Engagement}
	\end{center}
	\vspace{-4pt}
	\begin{itemize}[noitemsep, topsep=2pt, partopsep=0pt, parsep=0pt, leftmargin=*]
		\item \textbf{Selected Awards:} Synthetic Biology Challenges Silver Award $\times$ 2 (2025); Contemporary Undergraduate Mathematical Contest in Modeling National Second Prize (2025); Mathematical Contest in Modeling Honorable Mention (2025).
		\item \textbf{Academic Engagement:} Tsinghua-Peking University Center for Life Sciences Summer Camp (2025); Brown University Department of Physics Artificial Intelligence Winter School (2025).
	\end{itemize}
	
\end{document}